\chapter{ผลการทดลองและระบบต้นแบบ}
\label{chapter:result}

บทนี้รายงานผลการทดลองตามระเบียบวิธีในบทที่ \ref{chapter:experiment} แบ่งเป็นสามส่วน ได้แก่ ผลการทดลองหลักบนชุดข้อมูล CICIDS2017 แบบ 9 คลาส ระบบต้นแบบ และผลการทดลองเสริมเรื่องการส่งผ่านข้ามชุดข้อมูล ส่วนท้ายของบทรวบรวมข้อค้นพบสำคัญพร้อมตัวเลขไว้ในหน้าเดียว

%======================================================================
\section{ผลการทดลองหลักบนชุดข้อมูล CICIDS2017}
\label{sec:main-results}

\subsection{ประสิทธิภาพการจำแนกประเภท}
\label{sec:main-classification}

ตารางที่ \ref{tab:main-classification} แสดงผลของแบบจำลองทั้ง 7 รูปแบบบนชุดทดสอบ 749,401 ระเบียน ซึ่งแบ่งตามลำดับการจับข้อมูล (capture-order) แบบ 70/30 ภายในทุกกลุ่ม (ไฟล์, คลาส) คอลัมน์ Macro F1 (7 คลาส) คิดเฉพาะคลาสที่มีตัวอย่างในชุดทดสอบไม่น้อยกว่า 30 ระเบียน คือไม่รวม Heartbleed และ Infiltration

ค่า Macro F1-score ที่ใช้เป็นผลหลักและเกณฑ์จัดอันดับในงานจำแนก 9 คลาส คือค่าเฉลี่ย F1 ของทั้ง 9 คลาสโดยให้น้ำหนักแต่ละคลาสเท่ากัน ดังนั้นค่า 0.8873 ของ CatBoost เป็นค่าของทั้ง 9 คลาส ส่วนค่า 0.8613 ในคอลัมน์ 7 คลาสเป็นตัวชี้วัดประกอบเพื่อตรวจความไวของคะแนนต่อคลาสที่มีตัวอย่างน้อย โดยตัดเฉพาะ Heartbleed (ชุดทดสอบ 4 ระเบียน) และ Infiltration (11 ระเบียน) ออกจากการเฉลี่ย ไม่ได้ฝึกแบบจำลองใหม่หรือลบระเบียนเหล่านี้ออกจากชุดทดสอบ ผลรายคลาสยังรายงานครบทั้ง 9 คลาส และคะแนนจากคลาสที่มีตัวอย่างน้อยต้องตีความอย่างระมัดระวัง

ตัวชี้วัด 7 คลาสข้างต้นแตกต่างจาก Macro F1 ของการทดลองเสริมข้ามชุดข้อมูลในหัวข้อ \ref{sec:crossdataset} ซึ่งฝึกและประเมินบนโครงร่างคลาสร่วม 7 คลาส ได้แก่ Benign, DDoS, DoS, Brute Force, Bot, Infiltration และ Web Attack โดยตัด PortScan และ Heartbleed ที่ไม่มีคลาสคู่กันใน CSE-CIC-IDS2018 ออก จึงไม่ควรเปรียบเทียบค่าของทั้งสองกรณีโดยตรง แม้มีจำนวนคลาสเท่ากัน

% แหล่งข้อมูล: results/cicids2017_temporal_v1/report.md
\begin{table}[ht]
    \caption{ประสิทธิภาพการจำแนกประเภทของแบบจำลองทั้ง 7 รูปแบบบนชุดทดสอบของ CICIDS2017 (9 คลาส)}
    \label{tab:main-classification}
    \small
    \begin{tabularx}{\linewidth}{@{}l *{6}{>{\raggedleft\arraybackslash}X}@{}}
        \hline
        \textbf{แบบจำลอง} & \textbf{Accuracy} & \textbf{Balanced Acc.} & \textbf{Macro F1 (9 คลาส)} & \textbf{Macro F1 (7 คลาส)} & \textbf{Weighted F1} & \textbf{MCC} \\
        \hline
        CatBoost & 0.9953 & 0.8970 & \textbf{0.8873} & 0.8613 & 0.9952 & 0.9844 \\
        LightGBM & 0.9903 & 0.8576 & 0.8774 & 0.8627 & 0.9898 & 0.9676 \\
        XGBoost & 0.9906 & 0.8889 & 0.8773 & 0.8542 & 0.9902 & 0.9687 \\
        Random Forest & 0.9820 & 0.8165 & 0.8404 & 0.8378 & 0.9832 & 0.9397 \\
        Stacking Ensemble & 0.9581 & 0.8852 & 0.6822 & 0.7092 & 0.9712 & 0.8757 \\
        Logistic Regression & 0.8428 & 0.8338 & 0.5574 & 0.6348 & 0.9030 & 0.6651 \\
        Multi-Layer Perceptron & 0.9331 & 0.7112 & 0.5369 & 0.6767 & 0.9443 & 0.7895 \\
        \hline
    \end{tabularx}
\end{table}

CatBoost ได้ค่า Macro F1-score สูงสุดที่ 0.8873 ตามด้วย LightGBM (0.8774) และ XGBoost (0.8773) ซึ่งต่างกันไม่ถึง 0.01 แบบจำลองทุกตัวมีค่า Accuracy สูงกว่าแบบจำลองฐานที่ทำนาย BENIGN ทุกกรณี (0.8296) แต่ Logistic Regression สูงกว่าเพียงเล็กน้อย (0.8428) ซึ่งยืนยันว่าค่า Accuracy เพียงค่าเดียวไม่เพียงพอต่อการตัดสินบนข้อมูลที่ไม่สมดุล

รูปที่ \ref{fig:main-cm-catboost} แสดงผลการจำแนกของ CatBoost บนชุดทดสอบของรอบหลัก แถวแสดงคลาสจริงและคอลัมน์แสดงคลาสที่แบบจำลองทำนาย จาก Bot ทั้ง 585 ระเบียน แบบจำลองทำนายถูก 94 ระเบียน (ความไว 0.1607) และทำนายอีก 491 ระเบียนเป็น BENIGN ในอีกทิศทางหนึ่ง มี BENIGN 233 ระเบียนที่ถูกทำนายเป็น Bot สำหรับ DoS แบบจำลองทำนายผิดเป็น BENIGN 1,431 ระเบียน เป็น PortScan 913 ระเบียน และเป็น Web Attack 151 ระเบียน เมทริกซ์ของอีก 6 แบบจำลองแสดงในภาคผนวก \ref{appendix:main-confusion-matrices}

% แหล่งข้อมูล: results/cicids2017_temporal_v1/catboost_confusion_matrix.png
\begin{figure}[!htbp]
    \centering
    \includegraphics[width=0.92\linewidth]{catboost_confusion_matrix.png}
    \caption{เมทริกซ์ความสับสนของ CatBoost จากการทดลองหลักบนชุดทดสอบ CICIDS2017}
    \label{fig:main-cm-catboost}
\end{figure}
\FloatBarrier

\subsection{ความสามารถในการตรวจจับและผลกระทบเชิงปฏิบัติการ}

ตารางที่ \ref{tab:main-detection} ยุบผลการทำนายเป็นกรณีทวิภาคระหว่างการจราจรปกติและการโจมตี CatBoost ตรวจจับการโจมตีได้ 98.40\% โดยมีการแจ้งเตือนผิดพลาด 404 ครั้ง หรือ 6.5 ครั้งต่อการจราจรปกติหนึ่งหมื่นระเบียน ขณะที่ Logistic Regression มีอัตราการตรวจจับใกล้เคียงกัน (97.46\%) แต่แจ้งเตือนผิดพลาดถึง 105,258 ครั้ง

\begin{table}[ht]
    \caption{ความสามารถในการตรวจจับการโจมตีแบบทวิภาคบนชุดทดสอบของ CICIDS2017}
    \label{tab:main-detection}
    \small
    \begin{tabularx}{\linewidth}{@{}l *{5}{>{\raggedleft\arraybackslash}X}@{}}
        \hline
        \textbf{แบบจำลอง} & \textbf{อัตราตรวจจับ} & \textbf{FPR} & \textbf{แจ้งเตือนผิด} & \textbf{ต่อปกติ 10,000} & \textbf{การโจมตีที่พลาด} \\
        \hline
        CatBoost & 0.9840 & 0.00065 & 404 & 6.5 & 2,049 \\
        Logistic Regression & 0.9746 & 0.16931 & 105,258 & 1,693.1 & 3,244 \\
        Stacking Ensemble & 0.9661 & 0.04072 & 25,317 & 407.2 & 4,329 \\
        XGBoost & 0.9571 & 0.00038 & 236 & 3.8 & 5,483 \\
        LightGBM & 0.9513 & 0.00022 & 136 & 2.2 & 6,220 \\
        Random Forest & 0.9277 & 0.00530 & 3,298 & 53.0 & 9,232 \\
        Multi-Layer Perceptron & 0.8459 & 0.04721 & 29,347 & 472.1 & 19,677 \\
        \hline
    \end{tabularx}
\end{table}

เมื่อพิจารณารายคลาส คลาส Bot เป็นคลาสที่ตรวจจับได้แย่ที่สุดในทุกแบบจำลอง โดยมีค่าความไวเพียง 0.15--0.16 ส่วนคลาส DDoS และ PortScan มีค่าความไวเกิน 0.99 และ Brute Force เกิน 0.95 ในทุกแบบจำลอง คลาส Heartbleed (4 ระเบียน) และ Infiltration (11 ระเบียน) มีตัวอย่างในชุดทดสอบน้อยเกินกว่าจะตีความค่าความไวเป็นค่าที่เสถียรได้

\subsection{ประสิทธิภาพเชิงคำนวณและการจัดอันดับ}

ตารางที่ \ref{tab:main-efficiency} แสดงเวลาแฝงของการทำนายบนหน่วยประมวลผลกลาง โดยวัดเป็นชุดละ 1,000 โฟลว์ CatBoost มีเวลาแฝงที่ควอนไทล์ที่ 95 เท่ากับ 8.39 มิลลิวินาที ต่ำเป็นอันดับสองรองจาก Logistic Regression และมีขนาดไฟล์แบบจำลองเพียง 2.5 เมกะไบต์

\begin{table}[ht]
    \caption{ประสิทธิภาพเชิงคำนวณของแบบจำลอง (การทำนายบนหน่วยประมวลผลกลาง ชุดละ 1,000 โฟลว์)}
    \label{tab:main-efficiency}
    \small
    \begin{tabularx}{\linewidth}{@{}l *{5}{>{\raggedleft\arraybackslash}X}@{}}
        \hline
        \textbf{แบบจำลอง} & \textbf{เวลาฝึก (วินาที)} & \textbf{p50 (มิลลิวินาที)} & \textbf{p95 (มิลลิวินาที)} & \textbf{โฟลว์ต่อวินาที} & \textbf{ขนาด (MB)} \\
        \hline
        Logistic Regression & 653.7 & 4.30 & 5.15 & 232,391 & 0.0 \\
        CatBoost & 44.6 & 7.58 & 8.39 & 131,975 & 2.5 \\
        Multi-Layer Perceptron & 818.7 & 10.05 & 12.50 & 99,545 & 0.7 \\
        XGBoost & 154.6 & 43.42 & 54.10 & 23,030 & 15.8 \\
        Random Forest & 314.6 & 85.94 & 104.12 & 11,636 & 124.5 \\
        Stacking Ensemble & 2,754.7 & 119.79 & 137.95 & 8,348 & 109.3 \\
        LightGBM & 549.6 & 111.40 & 144.02 & 8,977 & 54.9 \\
        \hline
    \end{tabularx}
\end{table}

เมื่อนำเกณฑ์การจัดอันดับสามชุดที่ระบุไว้ก่อนการวิเคราะห์ผลมาใช้ CatBoost เป็นอันดับหนึ่งในทั้งสามเกณฑ์ ได้แก่ อันดับภาพรวม (Macro F1-score 0.8873) อันดับด้านความมั่นคงปลอดภัย (ความไวเฉลี่ยแบบ macro 0.8970 ภายใต้ FPR แบบทวิภาคไม่เกิน 0.01) และอันดับด้านการนำไปใช้งาน (p95 8.39 มิลลิวินาที ในกลุ่มแบบจำลองที่ Macro F1-score ต่างจากอันดับหนึ่งไม่เกิน 0.02) CatBoost จึงถูกเลือกเป็นแบบจำลองแชมเปียนของโครงงาน

\subsection{ผลการอธิบาย CatBoost ที่ฝึกซ้ำด้วย SHAP}

การวิเคราะห์ด้วย SHAP ตามวิธีในหัวข้อ \ref{sec:shap} ใช้ CatBoost ที่ฝึกซ้ำด้วยค่าไฮเปอร์พารามิเตอร์ในตารางที่ \ref{tab:best-params} บนชุดฝึกเดิมทั้ง 1,748,579 ระเบียน แบบจำลองที่ฝึกซ้ำได้ค่า Macro F1-score 0.8858 และ FPR 0.00063 บนชุดทดสอบเดิม ส่วน CatBoost ของรอบหลักที่ใช้สร้างเมทริกซ์ความสับสนในรูปที่ \ref{fig:main-cm-catboost} ได้ค่า Macro F1-score 0.8873 และ FPR 0.00065 ตามตารางที่ \ref{tab:main-classification} และ \ref{tab:main-detection} ค่า SHAP คำนวณด้วย TreeSHAP บนระเบียนจากชุดทดสอบที่สุ่มแยกตามคลาส คลาสละไม่เกิน 200 ระเบียน (Heartbleed มี 4 ระเบียนและ Infiltration มี 11 ระเบียนตามจำนวนที่มีในชุดทดสอบ) เพื่อไม่ให้ผลถูกครอบงำด้วยคลาส Benign ซึ่งมีสัดส่วนมากกว่า 80\% ของชุดทดสอบ

\begin{figure}[ht]
    \centering
    \includegraphics[width=0.8\linewidth]{fig_shap_main_bar.png}
    \caption{คุณลักษณะ 20 อันดับแรกของแบบจำลอง CatBoost เรียงตามค่าเฉลี่ยของค่าสัมบูรณ์ของค่า SHAP รวมทุกคลาส}
    \label{fig:shap-main-bar}
\end{figure}

รูปที่ \ref{fig:shap-main-bar} แสดงว่าในภาพรวม คุณลักษณะที่มีอิทธิพลสูงสุดคือ Destination Port (17.9\% ของผลรวมค่า SHAP) ตามด้วยขนาดหน้าต่าง TCP เริ่มต้นของทั้งสองทิศ คือ Init\_Win\_bytes\_forward (9.8\%) และ Init\_Win\_bytes\_backward (9.5\%) จากนั้นจึงเป็น min\_seg\_size\_forward (5.0\%) Flow Bytes/s (4.6\%) และกลุ่มช่วงเวลาระหว่างแพ็กเก็ตต่ำสุด ได้แก่ Flow IAT Min (4.3\%) และ Fwd IAT Min (3.5\%) ตารางที่ \ref{tab:shap-main-class} แสดงคุณลักษณะ 3 อันดับแรกของแต่ละคลาส

\begin{table}[ht]
    \caption{คุณลักษณะ 3 อันดับแรกของแต่ละคลาสจากการวิเคราะห์ SHAP ของแบบจำลอง CatBoost (ร้อยละของผลรวมค่า SHAP ของคลาสนั้น)}
    \label{tab:shap-main-class}
    \footnotesize
    \begin{tabularx}{\linewidth}{@{}l X X X@{}}
        \hline
        \textbf{คลาส} & \textbf{อันดับ 1} & \textbf{อันดับ 2} & \textbf{อันดับ 3} \\
        \hline
        Benign & Destination Port (16.3) & Init\_Win\_bytes\_forward (14.8) & Init\_Win\_bytes\_backward (9.6) \\
        DoS & Init\_Win\_bytes\_forward (18.3) & Destination Port (17.5) & Init\_Win\_bytes\_backward (6.4) \\
        DDoS & Init\_Win\_bytes\_backward (12.0) & Init\_Win\_bytes\_forward (10.6) & min\_seg\_size\_forward (10.3) \\
        PortScan & Init\_Win\_bytes\_forward (15.1) & Init\_Win\_bytes\_backward (10.7) & Destination Port (10.2) \\
        Brute Force & Destination Port (44.9) & Init\_Win\_bytes\_forward (6.2) & Fwd IAT Min (5.5) \\
        Bot & Destination Port (42.5) & Flow IAT Min (5.9) & Init\_Win\_bytes\_backward (5.8) \\
        Web Attack & Total Length of Fwd Packets (11.3) & Fwd Packet Length Max (10.8) & Fwd IAT Mean (6.3) \\
        Infiltration\textsuperscript{*} & Destination Port (26.1) & Init\_Win\_bytes\_forward (9.2) & min\_seg\_size\_forward (6.7) \\
        Heartbleed\textsuperscript{*} & Destination Port (25.9) & Init\_Win\_bytes\_backward (10.0) & Init\_Win\_bytes\_forward (6.8) \\
        \hline
        \multicolumn{4}{@{}l}{\scriptsize \textsuperscript{*} คำนวณจากระเบียนทดสอบเพียง 11 และ 4 ระเบียนตามลำดับ จึงใช้ดูแนวโน้มเท่านั้น} \\
    \end{tabularx}
\end{table}

ในระดับรายตัวอย่าง รูปที่ \ref{fig:shap-main-dos} แสดงค่า SHAP ของแต่ละระเบียนต่อผลลัพธ์คลาส DoS แต่ละจุดคือหนึ่งระเบียน ตำแหน่งตามแนวนอนคือขนาดและทิศทางที่คุณลักษณะนั้นผลักผลการทำนายเข้าหาหรือออกจาก DoS ส่วนสีแสดงค่าของคุณลักษณะ (น้ำเงินคือค่าต่ำ แดงคือค่าสูง) ระเบียนที่มีค่า Init\_Win\_bytes\_forward ต่ำถูกผลักไปทาง DoS อย่างชัดเจน (ค่า SHAP ประมาณ +2.5 ถึง +6) ขณะที่ระเบียนที่มีค่าสูงถูกผลักออก (ประมาณ −3) ส่วนค่า Destination Port ที่ต่ำ เช่น พอร์ต 80 ของเว็บเซิร์ฟเวอร์ ก็ผลักไปทาง DoS เช่นกัน

\begin{figure}[ht]
    \centering
    \includegraphics[width=0.72\linewidth]{fig_shap_main_dos.png}
    \caption{ค่า SHAP รายระเบียนของแบบจำลอง CatBoost ต่อผลลัพธ์คลาส DoS (15 คุณลักษณะแรก)}
    \label{fig:shap-main-dos}
\end{figure}

ข้อสังเกตจากผลการวิเคราะห์มีดังนี้

\begin{itemize}
    \item \textbf{แบบจำลองพึ่งพาหมายเลขพอร์ตปลายทางมาก} โดยเฉพาะ Brute Force และ Bot ที่ Destination Port คิดเป็นสัดส่วนเกือบครึ่งหนึ่งของค่า SHAP ทั้งหมด ซึ่งสอดคล้องกับลักษณะของการโจมตีที่มุ่งไปยังบริการเฉพาะ เช่น FTP (พอร์ต 21) และ SSH (พอร์ต 22) สำหรับ Brute Force และพอร์ต 8080 ซึ่ง Bot ใน CICIDS2017 มากกว่าครึ่งหนึ่งใช้ติดต่อกับเครื่องสั่งการ
    \item \textbf{คุณลักษณะระดับโพรโทคอล TCP มีบทบาทสูงในคลาสที่มีปริมาณมาก} ขนาดหน้าต่าง TCP เริ่มต้นและขนาดเซกเมนต์ต่ำสุดเป็นคุณลักษณะหลักของ DoS, DDoS และ PortScan ค่าเหล่านี้ขึ้นกับระบบปฏิบัติการและเครื่องมือที่ใช้สร้างทราฟฟิก มากกว่าพฤติกรรมของการโจมตีโดยตรง
    \item \textbf{Web Attack เป็นคลาสเดียวที่อาศัยคุณลักษณะด้านเนื้อหาและจังหวะเวลา} ได้แก่ ความยาวของข้อมูลขาไปและช่วงเวลาระหว่างแพ็กเก็ต ซึ่งสมเหตุสมผล เพราะการโจมตีเว็บใช้พอร์ต 80 เช่นเดียวกับทราฟฟิกปกติ
    \item \textbf{คุณลักษณะอันดับต้นซ้อนทับกับกลุ่มที่บ่งชี้ชุดข้อมูล} คุณลักษณะ 10 อันดับแรกในภาพรวมมี 8 ตัวที่อยู่ในกลุ่มคุณลักษณะที่บ่งชี้คอร์ปัส 13 ตัวของการทดลองเสริม (หัวข้อ \ref{sec:cd-shap-method}) แม้ทั้งสองการวิเคราะห์จะใช้โครงร่างคลาสและแบบจำลองต่างกัน ผลนี้ชี้ว่าคุณลักษณะที่แบบจำลองใช้จำแนกการโจมตีได้ดีภายใน CICIDS2017 เป็นกลุ่มเดียวกับคุณลักษณะที่มีค่าต่างกันมากระหว่างชุดข้อมูล ซึ่งเป็นคำอธิบายหนึ่งของผลการส่งผ่านที่จำกัดในหัวข้อ \ref{sec:cd-results}
\end{itemize}

%======================================================================
\FloatBarrier
\section{ระบบต้นแบบ}
\label{sec:prototype}

ระบบต้นแบบประกอบด้วยส่วนบริการเบื้องหลังที่พัฒนาด้วย FastAPI และส่วนติดต่อผู้ใช้ที่พัฒนาด้วย Next.js ส่วนติดต่อผู้ใช้แบ่งเป็นหน้าหลักดังนี้

\begin{itemize}
    \item \textbf{ภาพรวม} สรุปแบบจำลองแชมเปียนและตัวชี้วัดหลักของชุดผลลัพธ์ที่เลือก
    \item \textbf{ชุดข้อมูลและการสำรวจข้อมูล} แสดงการกระจายตัวของคลาสและผลการสำรวจข้อมูลเบื้องต้น
    \item \textbf{ประสิทธิภาพและการเปรียบเทียบ} แสดงตัวชี้วัดทั้งสี่มิติ เมทริกซ์ความสับสน และการจัดอันดับตามเกณฑ์ทั้งสามชุด
    \item \textbf{การอธิบายแบบจำลอง} แสดงคุณลักษณะที่มีอิทธิพลสูงสุดจากการวิเคราะห์ SHAP ทั้งภาพรวมและรายคลาส พร้อมค่าความเสถียรระหว่างรอบการสุ่ม
    \item \textbf{การทำนาย} รับไฟล์ข้อมูลโฟลว์ใหม่ ตรวจสอบโครงสร้างคอลัมน์เทียบกับรายชื่อคุณลักษณะที่บันทึกไว้พร้อมแบบจำลอง แล้วทำนายผลเป็นชุด
\end{itemize}

ภาพหน้าจอของแต่ละหน้าแสดงในรูปที่ \ref{fig:dash-overview} ถึง \ref{fig:dash-predict} ทุกภาพแสดงชุดผลลัพธ์ \texttt{cicids2017\_temporal\_v1} ของการทดลองหลัก (ชื่อโฟลเดอร์ในโค้ด ซึ่งใช้การแบ่งตามลำดับการจับข้อมูลตามหัวข้อ \ref{sec:split})

\begin{figure}[ht]
    \centering
    \includegraphics[width=\linewidth]{fig_dash_overview.png}
    \caption{หน้าภาพรวมของเว็บแดชบอร์ด แสดงแบบจำลองแชมเปียน ตัวชี้วัดหลัก และรายชื่อแบบจำลองในชุดผลลัพธ์}
    \label{fig:dash-overview}
\end{figure}

\begin{figure}[ht]
    \centering
    \includegraphics[width=\linewidth]{fig_dash_dataset.png}
    \caption{หน้าชุดข้อมูล แสดงจำนวนระเบียนและการกระจายตัวของคลาสในชุดฝึกและชุดทดสอบ}
    \label{fig:dash-dataset}
\end{figure}

\begin{figure}[ht]
    \centering
    \includegraphics[width=\linewidth]{fig_dash_performance.png}
    \caption{หน้าประสิทธิภาพของแบบจำลอง แสดงตัวชี้วัดรวมและค่าตัวชี้วัดรายคลาสของแบบจำลองที่เลือก}
    \label{fig:dash-performance}
\end{figure}

\begin{figure}[ht]
    \centering
    \includegraphics[width=\linewidth]{fig_dash_compare.png}
    \caption{หน้าเปรียบเทียบ แสดงตัวชี้วัดของแบบจำลองทั้ง 7 รูปแบบในตารางเดียว}
    \label{fig:dash-compare}
\end{figure}

\begin{figure}[ht]
    \centering
    \includegraphics[width=\linewidth]{fig_dash_predict.png}
    \caption{หน้าการทำนาย สำหรับเลือกแบบจำลองและอัปโหลดไฟล์ CSV ของโฟลว์ใหม่เพื่อทำนายผลเป็นชุด}
    \label{fig:dash-predict}
\end{figure}

%======================================================================
\FloatBarrier
\section{ผลการทดลองเสริม: การส่งผ่านข้ามชุดข้อมูล}
\label{sec:cd-results}

ผลในหัวข้อนี้ทั้งหมดมาจาก protocol\_v3 (ชุดฝึก 70\% / ชุดทดสอบ 30\%) ตามที่อธิบายในหัวข้อ \ref{sec:crossdataset} ผลของ protocol\_v2 (ชุดฝึก 60\%) ถูกอ้างเฉพาะในการวิเคราะห์ความไวต่อสัดส่วนการแบ่ง (หัวข้อ \ref{sec:cd-ssh})

\subsection{ผลหลัก}

ตารางที่ \ref{tab:cd-main-random} และ \ref{tab:cd-main-capture} แสดงค่า Macro F1-score ของการส่งผ่าน (ค่าเฉลี่ย ± ส่วนเบี่ยงเบนมาตรฐานจาก 5 รอบ) ค่า Recovery และค่าเพดานของแต่ละแบบจำลอง ค่าเพดานในแต่ละทิศคือค่าที่แบบจำลองชนิดเดียวกันทำได้ภายในคอร์ปัสปลายทาง ค่าพื้นของทิศ 2017\dirarrow{}2018 คือ 0.1337 และของทิศ 2018\dirarrow{}2017 คือ 0.1294 ตัวหนาคือค่า Macro F1-score สูงสุดของแต่ละทิศ

% แหล่งข้อมูล: docs/experiments/crossdataset/round6_split_70_30/data/summary_macro_f1.csv และ summary_gap.csv
\begin{table}[ht]
    \caption{ผลการส่งผ่านข้ามชุดข้อมูลภายใต้การแบ่งแบบสุ่ม (protocol\_v3)}
    \label{tab:cd-main-random}
    \footnotesize
    \begin{tabularx}{\linewidth}{@{}l c *{2}{>{\raggedleft\arraybackslash}X} c *{2}{>{\raggedleft\arraybackslash}X}@{}}
        \hline
        & \multicolumn{3}{c}{\textbf{2017\dirarrow{}2018 (ค่าพื้น 0.1337)}} & \multicolumn{3}{c}{\textbf{2018\dirarrow{}2017 (ค่าพื้น 0.1294)}} \\
        \textbf{แบบจำลอง} & \textbf{Macro F1} & \textbf{Rec.\ (\%)} & \textbf{เพดาน} & \textbf{Macro F1} & \textbf{Rec.\ (\%)} & \textbf{เพดาน} \\
        \hline
        LightGBM & \textbf{0.2850} ± 0.0181 & 20.6 & 0.8696 & 0.2860 ± 0.0524 & 18.4 & 0.9795 \\
        XGBoost & 0.2727 ± 0.0633 & 18.9 & 0.8676 & 0.2446 ± 0.0179 & 13.7 & 0.9715 \\
        CatBoost & 0.1692 ± 0.0073 & 4.9 & 0.8547 & 0.2195 ± 0.0176 & 10.7 & 0.9686 \\
        Random Forest & 0.1743 ± 0.0086 & 5.6 & 0.8625 & 0.1492 ± 0.0024 & 2.4 & 0.9612 \\
        Stacking & 0.2836 ± 0.0094 & 21.3 & 0.8383 & \textbf{0.3187} ± 0.0296 & 22.4 & 0.9727 \\
        MLP & 0.1489 ± 0.0068 & 2.2 & 0.8326 & 0.2256 ± 0.0173 & 13.0 & 0.8679 \\
        Logistic Regression & 0.1381 ± 0.0007 & 0.7 & 0.7811 & 0.2556 ± 0.0271 & 25.5 & 0.6243 \\
        \hline
        \textbf{เฉลี่ย 7 แบบจำลอง} & & 10.6 & & & 15.2 & \\
        \hline
    \end{tabularx}
\end{table}

\begin{table}[ht]
    \caption{ผลการส่งผ่านข้ามชุดข้อมูลภายใต้การแบ่งตามลำดับการจับข้อมูล (protocol\_v3)}
    \label{tab:cd-main-capture}
    \footnotesize
    \begin{tabularx}{\linewidth}{@{}l c *{2}{>{\raggedleft\arraybackslash}X} c *{2}{>{\raggedleft\arraybackslash}X}@{}}
        \hline
        & \multicolumn{3}{c}{\textbf{2017\dirarrow{}2018 (ค่าพื้น 0.1337)}} & \multicolumn{3}{c}{\textbf{2018\dirarrow{}2017 (ค่าพื้น 0.1294)}} \\
        \textbf{แบบจำลอง} & \textbf{Macro F1} & \textbf{Rec.\ (\%)} & \textbf{เพดาน} & \textbf{Macro F1} & \textbf{Rec.\ (\%)} & \textbf{เพดาน} \\
        \hline
        LightGBM & 0.3132 ± 0.0569 & 24.5 & 0.8675 & \textbf{0.4435} ± 0.0041 & 42.0 & 0.8767 \\
        XGBoost & \textbf{0.3368} ± 0.0599 & 28.0 & 0.8589 & 0.3524 ± 0.0187 & 30.4 & 0.8617 \\
        CatBoost & 0.1735 ± 0.0017 & 5.5 & 0.8516 & 0.2176 ± 0.0042 & 11.9 & 0.8711 \\
        Random Forest & 0.1717 ± 0.0001 & 5.3 & 0.8466 & 0.1631 ± 0.0056 & 4.6 & 0.8550 \\
        Stacking & 0.3322 ± 0.0600 & 28.6 & 0.8272 & 0.3567 ± 0.0217 & 30.9 & 0.8648 \\
        MLP & 0.1753 ± 0.0015 & 6.1 & 0.8145 & 0.3005 ± 0.0219 & 27.5 & 0.7511 \\
        Logistic Regression & 0.1496 ± 0.0000 & 2.6 & 0.7534 & 0.3506 ± 0.0000 & 50.1 & 0.5706 \\
        \hline
        \textbf{เฉลี่ย 7 แบบจำลอง} & & 14.4 & & & 28.2 & \\
        \hline
    \end{tabularx}
\end{table}

\begin{figure}[ht]
    \centering
    \includegraphics[width=\linewidth]{fig_cd_split_modes.png}
    \caption{ค่า Macro F1-score ของการส่งผ่านภายใต้การแบ่งแบบสุ่มและการแบ่งตามลำดับการจับข้อมูล (เส้นประคือค่าพื้นของชุดทดสอบปลายทาง)}
    \label{fig:cd-split-modes}
\end{figure}

ข้อสังเกตจากตารางที่ \ref{tab:cd-main-random}, \ref{tab:cd-main-capture} และรูปที่ \ref{fig:cd-split-modes} มีดังนี้

\begin{itemize}
    \item \textbf{แบบจำลองที่ได้ค่าสูงสุดในแต่ละทิศ} ภายใต้การแบ่งตามลำดับการจับข้อมูล ทิศ 2017\dirarrow{}2018 คือ XGBoost (0.3368) ทิศ 2018\dirarrow{}2017 คือ LightGBM (0.4435) ภายใต้การแบ่งแบบสุ่ม ทิศ 2017\dirarrow{}2018 คือ LightGBM (0.2850) ทิศ 2018\dirarrow{}2017 คือ Stacking (0.3187) ค่าสูงสุดทุกกรณียังห่างจากค่าเพดาน (0.57--0.98) มาก
    \item \textbf{การแบ่งแบบสุ่มกับการแบ่งตามลำดับการจับข้อมูล} ค่าเพดานของ CICIDS2017 ภายใต้การแบ่งแบบสุ่มสูงกว่าการแบ่งตามลำดับการจับข้อมูลเฉลี่ย 0.099 (0.9065 เทียบกับ 0.8073) ขณะที่ของ CSE-CIC-IDS2018 ต่างกันเพียง 0.012 (0.8438 เทียบกับ 0.8314) ผลนี้แสดงความแตกต่างของค่าเพดานระหว่างสองวิธีแบ่งข้อมูล โดยยังไม่สามารถตัดสินว่าค่าใดสะท้อนประสิทธิภาพการใช้งานจริงได้ดีกว่า เนื่องจากการแบ่งตามลำดับการจับข้อมูลไม่ใช่การกันข้อมูลตามเวลาประทับ และองค์ประกอบของเครื่องมือโจมตีย่อยมีผลต่อคะแนนด้วย ค่า Recovery จึงต้องอ่านคู่กับค่าดิบและค่าเพดานของแต่ละวิธีแบ่งข้อมูล
    \item \textbf{ความไม่สมมาตร} ภายใต้การแบ่งตามลำดับการจับข้อมูล ค่า Recovery เฉลี่ยของทิศ 2018\dirarrow{}2017 (28.2\%) สูงกว่าทิศ 2017\dirarrow{}2018 (14.4\%) อยู่ 13.8 จุด และทิศ 2018\dirarrow{}2017 ดีกว่าใน 6 จาก 7 แบบจำลอง (Random Forest ต่างกันไม่ถึง 1 จุด) ภายใต้การแบ่งแบบสุ่ม ค่า Recovery เฉลี่ยของทิศ 2018\dirarrow{}2017 และ 2017\dirarrow{}2018 เท่ากับ 15.2\% และ 10.6\% ตามลำดับ ความไม่สมมาตรที่สังเกตได้จึงมีขนาดเล็กลงเป็น 4.6 จุดเปอร์เซ็นต์ และทิศ 2018\dirarrow{}2017 ดีกว่าใน 4 จาก 7 แบบจำลอง ความไม่สมมาตรจึงยังคงอยู่ แต่ขนาดขึ้นกับวิธีแบ่งข้อมูลและสัดส่วนการแบ่ง (หัวข้อ \ref{sec:cd-ssh})
    \item \textbf{การจัดอันดับต้องดูค่าดิบคู่กับค่า Recovery} Logistic Regression ได้ค่า Recovery สูงสุดในทิศ 2018\dirarrow{}2017 (50.1\%) แต่ค่า Macro F1-score ดิบ (0.3506) ต่ำกว่า LightGBM (0.4435) ที่ได้ Recovery 42.0\% เพราะค่าเพดานของ Logistic Regression บน CICIDS2017 ต่ำเพียง 0.5706 การใช้ค่า Recovery เพียงค่าเดียวจะจัดอันดับผิด
    \item \textbf{ระดับของค่าดิบ} ในทิศ 2017\dirarrow{}2018 มี 4 จาก 7 แบบจำลอง (CatBoost, Random Forest, MLP และ Logistic Regression) ที่ได้ค่า Macro F1-score เพียง 0.15--0.18 ใกล้กับค่าพื้น 0.1337
\end{itemize}

\subsection{ผลรายคลาส}
\label{sec:cd-per-class}

ตารางที่ \ref{tab:cd-per-class} และรูปที่ \ref{fig:cd-per-class} แสดงค่า F1 รายคลาสภายใต้การแบ่งตามลำดับการจับข้อมูล ค่าเพดานในตารางเป็นค่าเฉลี่ยของ 7 แบบจำลอง ส่วนในรูปเป็นค่าสูงสุดของ 7 แบบจำลอง

% แหล่งข้อมูล: round6_split_70_30/data/summary_per_class.csv และ report/protocol_v3_per_class.md
\begin{table}[ht]
    \caption{ค่า F1 รายคลาสของการส่งผ่านภายใต้การแบ่งตามลำดับการจับข้อมูล (protocol\_v3)}
    \label{tab:cd-per-class}
    \footnotesize
    \begin{tabularx}{\linewidth}{@{}l l l c c X@{}}
        \hline
        \textbf{คลาส} & \textbf{สูงสุด 2017\dirarrow{}2018} & \textbf{สูงสุด 2018\dirarrow{}2017} & \textbf{เพดาน 2017} & \textbf{เพดาน 2018} & \textbf{การส่งผ่าน} \\
        \hline
        Brute Force & 0.929 (Stacking) & 0.969 (XGBoost) & 0.990 & 0.999 & ดีในบางแบบจำลอง \\
        DoS & 0.279 (MLP) & 0.590 (CatBoost) & 0.910 & 0.939 & บางส่วน \\
        DDoS & 0.011 (MLP) & 0.518 (LR) & 0.993 & 0.991 & แย่ \\
        Bot & 0.358 (XGBoost) & 0.869 (LightGBM) & 0.208 & 0.998 & ขึ้นกับแบบจำลอง \\
        Web Attack & 0.002 (XGBoost) & 0.018 (LightGBM) & 0.803 & 0.782 & แย่ \\
        Benign & 0.942 (XGBoost) & 0.954 (LR) & 0.987 & 0.991 & ดี \\
        Infiltration & 0.000 & 0.000 & 0.760 & 0.120 & รายงานไม่ได้ \\
        \hline
    \end{tabularx}
\end{table}

\begin{figure}[ht]
    \centering
    \includegraphics[width=0.72\linewidth]{fig_cd_per_class.png}
    \caption{ค่า F1 สูงสุดจาก 7 แบบจำลองแยกตามคลาส ภายใต้การแบ่งตามลำดับการจับข้อมูล (\textsuperscript{*} Infiltration มีตัวอย่างทดสอบใน CICIDS2017 เพียง 11 ระเบียน)}
    \label{fig:cd-per-class}
\end{figure}

\begin{itemize}
    \item \textbf{Brute Force} ส่งผ่านได้ดีเมื่อแบบจำลองเป็น LightGBM, XGBoost หรือ Stacking ในทิศ 2017\dirarrow{}2018 (0.77--0.93) และในทิศ 2018\dirarrow{}2017 มี 5 แบบจำลองที่ได้ตั้งแต่ 0.5 ขึ้นไป แต่ CatBoost และ Random Forest ได้ไม่ถึง 0.1 ในทั้งสองทิศ
    \item \textbf{DoS} ได้ค่าต่ำถึงปานกลางในทั้งสองทิศ (0.12--0.28 และ 0.19--0.59) เทียบกับค่าเพดาน 0.79--0.98 และไม่เปลี่ยนเมื่อเปลี่ยนสัดส่วนการแบ่ง
    \item \textbf{DDoS} มีค่าเพดานเกือบสมบูรณ์ทั้งสองคอร์ปัส แต่ส่งผ่านในทิศ 2017\dirarrow{}2018 ได้สูงสุดเพียง 0.011 และในทิศ 2018\dirarrow{}2017 มีเพียง Logistic Regression ที่ได้เกิน 0.14 (0.518) DDoS จึงเป็นคลาสที่สูญเสียมากที่สุดเมื่อเทียบค่าเพดานกับค่าการส่งผ่าน
    \item \textbf{Bot} จำแนกได้ยากแม้ภายใน CICIDS2017 เอง (ค่าเพดาน 0.00--0.27) แต่ภายใน CSE-CIC-IDS2018 ได้เกือบ 1.0 ในทิศ 2018\dirarrow{}2017 มีเพียง LightGBM ที่ได้เกิน 0.05 (0.869)
    \item \textbf{Web Attack} ส่งผ่านไม่ได้ในทั้งสองทิศ ทั้งที่ค่าเพดานสูง
    \item \textbf{Benign} ได้ 0.88--0.95 ในทุกกรณี เนื่องจาก Macro F1-score เป็นค่าเฉลี่ยของ 7 คลาส คลาส Benign เพียงคลาสเดียวจึงคิดเป็นคะแนนประมาณ 0.13 ซึ่งเท่ากับค่าพื้นโดยประมาณ คะแนนส่วนที่สูงกว่าค่าพื้นจึงมาจากคลาสโจมตีเพียงไม่กี่คลาส ได้แก่ Brute Force, DoS และ Bot ในบางแบบจำลอง
    \item \textbf{Infiltration} รายงานไม่ได้อย่างน่าเชื่อถือ เพราะ CICIDS2017 มีตัวอย่างเพียง 36 ระเบียน (ชุดฝึก 25 ชุดทดสอบ 11) ค่า F1 ของทิศ 2018\dirarrow{}2017 จึงคำนวณจากระเบียนเพียง 11 ระเบียน ขณะที่แม้ภายใน CSE-CIC-IDS2018 เองแบบจำลองก็ได้ค่า F1 เพียง 0.00--0.22
\end{itemize}

\subsection{ความไวต่อสัดส่วนการแบ่งข้อมูล: กรณี SSH-Patator}
\label{sec:cd-ssh}

เมื่อเปลี่ยนสัดส่วนชุดฝึกจาก 60\% (protocol\_v2) เป็น 70\% (protocol\_v3) ผลภายใต้การแบ่งแบบสุ่มแทบไม่เปลี่ยน แต่ผลภายใต้การแบ่งตามลำดับการจับข้อมูลเปลี่ยนมาก สาเหตุอยู่ที่องค์ประกอบของคลาส Brute Force ใน CICIDS2017 ซึ่งประกอบด้วยเครื่องมือโจมตี 2 ตัวที่ทำงานต่อเนื่องกันในไฟล์เดียว คือ FTP-Patator แล้วตามด้วย SSH-Patator การแบ่งตามลำดับการจับข้อมูลตัดที่สัดส่วนของแถวในกลุ่ม ตำแหน่งที่ตัดจึงกำหนดว่าเครื่องมือใดอยู่ในชุดฝึก ดังตารางที่ \ref{tab:cd-ssh}

\begin{table}[ht]
    \caption{องค์ประกอบของคลาส Brute Force ใน CICIDS2017 ภายใต้การแบ่งตามลำดับการจับข้อมูล}
    \label{tab:cd-ssh}
    \small
    \begin{tabularx}{\linewidth}{@{}X r r r@{}}
        \hline
        \textbf{เครื่องมือโจมตี} & \textbf{ชุดฝึก (60\%)} & \textbf{ชุดฝึก (70\%)} & \textbf{ชุดทดสอบ (ทั้งสองแบบ)} \\
        \hline
        FTP-Patator & 5,490 & 5,931 & 0 \\
        SSH-Patator & \textbf{0} & \textbf{474} & 2,745 \\
        \hline
    \end{tabularx}
\end{table}

ที่ชุดฝึก 60\% แบบจำลองที่ฝึกจาก CICIDS2017 ไม่เคยเห็น SSH-Patator เลย ขณะที่ชุดทดสอบของคลาสนี้เป็น SSH-Patator ทั้งหมด 2,745 ระเบียน ที่ชุดฝึก 70\% รอยตัดเลื่อนเข้าไปในช่วงของ SSH-Patator ทำให้มี SSH-Patator 474 ระเบียนในชุดฝึก ซึ่งเป็นเครื่องมือเดียวกับ SSH-Bruteforce ของ CSE-CIC-IDS2018 ผลที่ตามมาแสดงในตารางที่ \ref{tab:cd-ssh-f1} และรูปที่ \ref{fig:cd-split-sensitivity}

\begin{table}[ht]
    \caption{ค่า F1 ของคลาส Brute Force และ Macro F1-score ในทิศ 2017\dirarrow{}2018 ภายใต้การแบ่งตามลำดับการจับข้อมูล}
    \label{tab:cd-ssh-f1}
    \small
    \begin{tabularx}{\linewidth}{@{}l *{4}{>{\raggedleft\arraybackslash}X}@{}}
        \hline
        \textbf{แบบจำลอง} & \textbf{Brute Force (60\%)} & \textbf{Brute Force (70\%)} & \textbf{Macro F1 (60\%)} & \textbf{Macro F1 (70\%)} \\
        \hline
        LightGBM & 0.0000 & 0.7729 & 0.2058 & 0.3132 \\
        XGBoost & 0.0000 & 0.7901 & 0.1712 & 0.3368 \\
        Stacking & 0.0000 & 0.9290 & 0.1715 & 0.3322 \\
        CatBoost / RF / MLP & 0.0000 & 0.0000 & 0.17--0.18 & 0.17--0.18 \\
        Logistic Regression & 0.0024 & 0.0000 & 0.1508 & 0.1496 \\
        \hline
    \end{tabularx}
\end{table}

\begin{figure}[ht]
    \centering
    \includegraphics[width=\linewidth]{fig_cd_split_sensitivity.png}
    \caption{ผลของการเปลี่ยนสัดส่วนชุดฝึกจาก 60\% เป็น 70\% ต่อค่า Macro F1-score (ซ้าย) และค่า F1 ของคลาส Brute Force (ขวา) ในทิศ 2017\dirarrow{}2018}
    \label{fig:cd-split-sensitivity}
\end{figure}

เนื่องจาก Macro F1-score เป็นค่าเฉลี่ยของ 7 คลาส การเปลี่ยนของ Macro F1-score จึงแยกได้เป็นผลรวมของการเปลี่ยนรายคลาส (คลาสละ ผลต่างของ F1 หารด้วย 7) ตารางที่ \ref{tab:cd-delta-class} แสดงว่าคลาสใดเป็นเจ้าของการเปลี่ยนแปลง

\begin{table}[ht]
    \caption{การเปลี่ยนของ Macro F1-score จากชุดฝึก 60\% เป็น 70\% แยกตามคลาส (การแบ่งตามลำดับการจับข้อมูล)}
    \label{tab:cd-delta-class}
    \footnotesize
    \begin{tabularx}{\linewidth}{@{}l *{6}{>{\raggedleft\arraybackslash}X}@{}}
        \hline
        \textbf{ค่าที่วัด} & \textbf{ผลต่าง Macro F1} & \textbf{Brute Force} & \textbf{Bot} & \textbf{DoS} & \textbf{DDoS} & \textbf{คลาสอื่น} \\
        \hline
        เพดาน 2017 (เฉลี่ย 7 แบบจำลอง) & +0.158 & +0.141 & 0.000 & 0.000 & 0.000 & +0.017 \\
        2017\dirarrow{}2018 (เฉลี่ย 7 แบบจำลอง) & +0.062 & +0.051 & +0.011 & 0.000 & 0.000 & 0.000 \\
        \quad LightGBM & +0.107 & +0.110 & 0.000 & -0.003 & 0.000 & 0.000 \\
        \quad XGBoost & +0.166 & +0.113 & +0.051 & +0.001 & 0.000 & +0.001 \\
        \quad Stacking & +0.161 & +0.133 & +0.027 & 0.000 & 0.000 & +0.001 \\
        2018\dirarrow{}2017 (เฉลี่ย 7 แบบจำลอง) & -0.006 & -0.001 & 0.000 & 0.000 & -0.004 & -0.001 \\
        \hline
    \end{tabularx}
\end{table}

\begin{itemize}
    \item ค่าเพดานของ CICIDS2017 ที่เพิ่มขึ้น 0.158 มาจาก Brute Force 0.141 (89\%) ส่วนที่เหลือเกือบทั้งหมดคือ Infiltration ซึ่งมีชุดทดสอบเพียง 11 ระเบียน
    \item ค่า Macro F1-score ของการส่งผ่านในทิศ 2017\dirarrow{}2018 ที่เพิ่มขึ้น 0.062 มาจาก Brute Force 82\% และ Bot 18\% การเปลี่ยนเกิดเฉพาะใน LightGBM, XGBoost และ Stacking ส่วนอีก 4 แบบจำลองเปลี่ยนไม่ถึง 0.002 ส่วน Bot ที่เพิ่มขึ้นใน XGBoost (F1 จาก 0.000 เป็น 0.358) และ Stacking (จาก 0.000 เป็น 0.187) ไม่ได้มาจากองค์ประกอบของชุดฝึก เมื่อแยกดูราย seed พบว่าค่า F1 ของ Bot ออกมาเพียงสองแบบ คือเกือบ 1.0 หรือ 0 ทั้งหมด XGBoost จำแนก Bot ได้ใน 0 จาก 5 รอบที่ชุดฝึก 60\% และ 2 จาก 5 รอบที่ชุดฝึก 70\% Stacking ได้ 0 และ 1 จาก 5 รอบ ส่วน LightGBM ได้ 1 จาก 5 รอบเท่ากันทั้งสองแบบ นอกจากนี้ Bot 195 ระเบียนที่เข้ามาในชุดฝึกเมื่อเลื่อนรอยตัดก็ไม่ได้ใกล้เคียง Bot ของ CSE-CIC-IDS2018 มากกว่า Bot ที่อยู่ในชุดฝึกอยู่แล้ว (มี Bot ของ CSE-CIC-IDS2018 เพียง 0.1\% ที่มีเพื่อนบ้านใกล้ที่สุดอยู่ในกลุ่มนี้ เทียบกับสัดส่วน 14.3\% ของกลุ่มนี้ในชุดฝึก) ค่าที่เพิ่มขึ้นจึงสะท้อนความไม่เสถียรของแบบจำลองระหว่างรอบการสุ่ม ไม่ใช่ผลของการเปลี่ยนสัดส่วนการแบ่ง
    \item ในทิศ 2018\dirarrow{}2017 ค่าเฉลี่ยแทบไม่เปลี่ยน แต่ใน Stacking คลาส DDoS ลดลงจาก 0.324 เป็น 0.132 ขณะที่ DoS เพิ่มจาก 0.384 เป็น 0.493 สองคลาสหักล้างกัน Macro F1-score จึงดูคงที่
    \item ความไม่สมมาตรภายใต้การแบ่งตามลำดับการจับข้อมูลลดลงจาก 33.6 จุด (ชุดฝึก 60\%) เป็น 13.8 จุด (ชุดฝึก 70\%) และค่าเพดานเฉลี่ยของ CICIDS2017 เพิ่มจาก 0.649 เป็น 0.807 ส่วนภายใต้การแบ่งแบบสุ่ม ความไม่สมมาตรเปลี่ยนจาก 4.3 เป็น 4.6 จุดเท่านั้น
\end{itemize}

ข้อสรุปของหัวข้อนี้คือ ภายใต้การแบ่งตามลำดับการจับข้อมูล องค์ประกอบของเครื่องมือโจมตีย่อยในชุดฝึกสามารถเปลี่ยนข้อสรุปเชิงปริมาณของการทดลองข้ามชุดข้อมูลได้มาก เพียงการเลื่อนรอยตัดให้ครอบคลุม SSH-Patator 474 ระเบียนทำให้ความไม่สมมาตรลดลงกว่าครึ่ง ขนาดของความไม่สมมาตรจึงควรรายงานเป็นช่วง (14--34 จุด) พร้อมระบุองค์ประกอบของชุดฝึก มิใช่ตัวเลขเดียว ขณะที่ทิศทางของความไม่สมมาตรยังคงเดิม

\subsection{ความแตกต่างเฉพาะชุดข้อมูล}
\label{sec:cd-av}

ผลการตรวจสอบเชิงปฏิปักษ์ตามวิธีในหัวข้อ \ref{sec:cd-av-method} แสดงในตารางที่ \ref{tab:cd-av} และรูปที่ \ref{fig:cd-av} การตรวจนี้ฝึกตัวจำแนกให้แยกคอร์ปัสบนแถวที่สุ่มจากแต่ละคลาสด้วยการแบ่งของตนเอง จึงไม่ขึ้นกับสัดส่วนการแบ่ง 60/10/30 หรือ 70/30 ของโพรโทคอลการส่งผ่าน

% แหล่งข้อมูล: round5_reliability_audit/data/diagnostics.json, within_dataset_control.json, tool_matched.json
\begin{table}[ht]
    \caption{ค่า AUC ของการแยกคอร์ปัส CICIDS2017 กับ CSE-CIC-IDS2018 รายคลาส (ค่าเฉลี่ย 5 รอบ)}
    \label{tab:cd-av}
    \footnotesize
    \begin{tabularx}{\linewidth}{@{}X r r r l r r@{}}
        \hline
        \textbf{คลาส} & \textbf{LightGBM} & \textbf{LR} & \textbf{Stump} & \textbf{ภายในคอร์ปัส} & \textbf{ซ้ำทุกค่า} & \textbf{ใกล้ซ้ำ} \\
        \hline
        Benign & 0.9990 & 0.9478 & 0.7089 & 0.6580--0.8112 & 0 & 0.10\% \\
        DDoS & 1.0000 & 1.0000 & 0.9917 & 0.9997--1.0000\textsuperscript{†} & 0 & 6.82\% \\
        DoS & 1.0000 & 0.9951 & 0.9370 & 0.9952--1.0000\textsuperscript{†} & 0 & 0.16\% \\
        Brute Force & 1.0000 & 0.9999 & 0.9950 & 0.8606--0.8993 & 0 & 0.20\% \\
        Bot & 1.0000 & 0.9990 & 0.9959 & 0.7584--0.9365 & 0 & 0.17\% \\
        Web Attack & 1.0000 & 1.0000 & 0.9992 & 0.7095--0.8404 & 0 & 2.75\% \\
        \hline
        \multicolumn{7}{@{}l}{\scriptsize ภายในคอร์ปัส: ช่วงของค่า AUC จากแบบจำลองทั้ง 7 รูปแบบ เมื่อแยกคนละไฟล์หรือคนละช่วงเวลาในคอร์ปัสเดียวกัน} \\
        \multicolumn{7}{@{}l}{\scriptsize \textsuperscript{†} ใน CSE-CIC-IDS2018 สองฝั่งของเส้นแบ่งเป็นเครื่องมือโจมตีคนละตัว ไม่ใช่เพียงคนละวัน} \\
    \end{tabularx}
\end{table}

\begin{figure}[ht]
    \centering
    \includegraphics[width=0.9\linewidth]{fig_cd_av.png}
    \caption{ช่วงของค่า AUC ในการแยกข้อมูลสองกลุ่ม ตั้งแต่กลุ่มควบคุมที่ไม่ควรแยกได้ไปจนถึงการแยกข้ามคอร์ปัส (กลุ่มควบคุมภายในคอร์ปัสไม่รวม DoS และ DDoS)}
    \label{fig:cd-av}
\end{figure}

\begin{itemize}
    \item การแยกข้ามคอร์ปัสได้ค่า AUC 0.9990--1.0000 ในทุกคลาสที่ทดสอบได้ และค่าแทบไม่ต่างกันระหว่าง 5 รอบ คลาส Infiltration ทดสอบไม่ได้ เพราะ CICIDS2017 มีเพียง 36 ระเบียน
    \item การสลับป้ายกำกับทำให้ค่า AUC ลดลงเหลือ 0.4928--0.5055 ในทุกแบบจำลองและทุกคลาส ส่วนกลุ่มควบคุมว่างได้ 0.49--0.53
    \item ภายในคอร์ปัสเดียวกัน คนละไฟล์หรือคนละช่วงเวลาแยกได้ 0.66--0.94 (ไม่รวม DoS และ DDoS) เมื่อตรึงเครื่องมือโจมตีให้เหมือนกันแล้วเปลี่ยนเฉพาะวัน ได้ 0.75--0.83 แต่เมื่อตรึงเครื่องมือโจมตีให้เหมือนกันแล้วเปลี่ยนเฉพาะคอร์ปัส ได้ 0.9999--1.0000
    \item Logistic Regression แยกได้ตั้งแต่ 0.9478 ขึ้นไป และ decision stump ซึ่งใช้กฎเดียวบนคอลัมน์เดียวแยกคลาสโจมตีได้ 0.937--0.999 (Benign ได้ 0.709)
    \item ไม่มีระเบียนที่ซ้ำกันทุกค่าคร่อมชุดฝึกและชุดทดสอบเลยในการสุ่มทั้ง 30 ครั้ง ระเบียนใกล้ซ้ำอยู่ใกล้ระดับบังเอิญ (0.10--0.20\%) ยกเว้น DDoS (6.82\%) และ Web Attack (2.75\%) แต่ decision stump ซึ่งจำแถวไม่ได้ก็ยังแยกสองคลาสนี้ได้ 0.99 ขึ้นไป
\end{itemize}

สรุปได้ว่าสองคอร์ปัสแยกออกจากกันได้เกือบสมบูรณ์ และความต่างนี้แรงกว่าความต่างระหว่างวันหรือไฟล์ภายในคอร์ปัสเดียวกันอย่างชัดเจน หลักฐานที่บ่งชี้ว่าค่า AUC ที่สูงไม่ได้มาจากการรั่วไหลแบบตรง ๆ มีสามข้อ คือ การสลับป้ายกำกับให้ค่าระดับสุ่ม ไม่มีระเบียนซ้ำคร่อมชุดฝึกและชุดทดสอบ และแบบจำลองที่มีความจุต่ำมากอย่าง decision stump ก็แยกได้ ความต่างดังกล่าวจึงอยู่ในค่าของคุณลักษณะเอง และกระจายอยู่ในคุณลักษณะจำนวนมาก โดยเฉพาะกลุ่มขนาดแพ็กเก็ต ซึ่งเป็นกลุ่มเดียวกับที่ใช้จำแนกการโจมตี (สหสัมพันธ์อันดับระหว่างความสามารถแยกคอร์ปัสกับ mutual information ต่อคลาสโจมตีเท่ากับ +0.72 ใน CICIDS2017 และ +0.77 ใน CSE-CIC-IDS2018)

\subsection{การวิเคราะห์สาเหตุของการทำนายพลาดด้วย SHAP}
\label{sec:cd-shap}

การทดสอบสมมติฐานตามหัวข้อ \ref{sec:cd-shap-method} บนการแบ่ง 70/30 ให้ผลดังนี้

\begin{itemize}
    \item การเทียบภายในคอร์ปัสปลายทาง (กลุ่ม F หลัก 13 ตัว) มีกรณีที่จำนวนข้อมูลพอให้ตัดสินได้ 11 กรณีจาก 4 แบบจำลองและสองทิศ (DoS 8 กรณี และอีก 3 กรณีคือ LightGBM Brute Force ทิศ 2017\dirarrow{}2018, LightGBM Bot และ XGBoost DDoS ทิศ 2018\dirarrow{}2017) \textbf{ไม่ผ่านเกณฑ์ทั้ง 11 กรณี} ค่า \textit{G}\textsubscript{within} ติดลบ 9 จาก 11 กรณี และค่า \textit{p} ต่ำสุดคือ 0.174
    \item กรณี DoS ทิศ 2017\dirarrow{}2018 ตอบคำถามได้ชัดที่สุด เมื่อเทียบกับระเบียนที่ทำนายถูกในคอร์ปัสต้นทาง ส่วนต่างเป็นบวกใน LightGBM, XGBoost และ CatBoost แต่เมื่อเทียบกับระเบียนที่ทำนายถูกในคอร์ปัสปลายทางเดียวกัน \textit{G}\textsubscript{within} ติดลบในทั้ง 4 แบบจำลอง (-0.07 ถึง -0.48) และติดลบทั้ง 20 ค่าราย seed แสดงว่าส่วนต่างบวกเดิมมาจากการที่ข้อมูลมาจากต่างคอร์ปัส ไม่ได้มาจากการทำนายผิด
    \item การเทียบกับกลุ่มคุณลักษณะสุ่มในทิศ 2017\dirarrow{}2018 (เทียบกับระเบียนที่ทำนายถูกในคอร์ปัสต้นทาง) มีกรณีที่นับได้ 17 กรณี ไม่ผ่านทั้งหมด ค่า \textit{p} ต่ำสุด 0.063 (CatBoost DoS)
\end{itemize}

ข้อค้นพบรองที่ต้องรายงานอย่างระมัดระวังคือ เมื่อใช้กลุ่มคุณลักษณะที่เลือกด้วยเกณฑ์ SHAP อย่างเดียว (26 ตัว) มี 2 กรณีที่ผ่านเกณฑ์รายคลาส ได้แก่ XGBoost DoS (\textit{G} = +0.228, \textit{p} = 0.0045) และ CatBoost DDoS (\textit{G} = +0.246, \textit{p} = 0.0090) ทั้งสองไม่ใช่กลุ่มหลักที่ใช้ตัดสิน และไม่ทำให้ผลรวมของแบบจำลองนั้นเปลี่ยน เมื่อตรวจซ้ำแบบภายในคอร์ปัสปลายทาง XGBoost DoS ได้ \textit{G}\textsubscript{within} = +0.073 (\textit{p} = 0.21) ซึ่งไม่ผ่าน ส่วน CatBoost DDoS ตรวจไม่ได้ เพราะแบบจำลองทำนาย DDoS ใน CSE-CIC-IDS2018 ถูก 0 ระเบียน นอกจากนี้ การวิเคราะห์รอบเสริมซึ่งระบุเกณฑ์ไว้ก่อนการวิเคราะห์เช่นกัน ได้ใช้กลุ่มคุณลักษณะ 13 ตัวที่ค่าของคลาสโจมตีนั้นเลื่อนมากที่สุดระหว่างสองคอร์ปัสแทนกลุ่ม F ผลรวมตามเกณฑ์เป็นแบบบางส่วน โดยคลาสที่ชัดที่สุดคือ DoS ซึ่งผ่านเกณฑ์ใน 3 จาก 4 แบบจำลอง อย่างไรก็ตาม รอบนี้เทียบกับระเบียนที่ทำนายถูกในคอร์ปัสต้นทาง ซึ่งการเทียบภายในคอร์ปัสปลายทางข้างต้นแสดงแล้วว่ายังปนผลของการที่ข้อมูลมาจากต่างคอร์ปัส ผลนี้จึงบ่งชี้ได้เพียงว่าการเลื่อนค่าระดับคลาสอาจมีส่วนต่อการทำนายพลาด และต้องตรวจซ้ำด้วยการเทียบภายในคอร์ปัสปลายทางก่อนนำไปสรุป

ข้อสรุปที่ใช้ได้อย่างปลอดภัยคือ \emph{ไม่พบหลักฐานว่ากลุ่มคุณลักษณะที่บ่งชี้ชุดข้อมูลซึ่งกำหนดไว้ก่อนการวิเคราะห์เป็นสาเหตุโดยตรงของการทำนายระเบียนโจมตีเป็น Benign หลังการส่งผ่านข้ามชุดข้อมูล}

%======================================================================
\clearpage
\section{สรุปข้อค้นพบสำคัญ}
\label{sec:key-findings}

ข้อค้นพบด้านล่างใช้ protocol\_v3 (ชุดฝึก 70\% / ชุดทดสอบ 30\%) เป็นหลัก ตัวเลขของชุดฝึก 60\% ใช้เพื่อเปรียบเทียบความไวเท่านั้น

\begin{enumerate}
    \item \textbf{การทดลองหลัก} CatBoost ได้ Macro F1-score สูงสุดบน CICIDS2017 แบบ 9 คลาสที่ 0.8873 ตรวจจับการโจมตีได้ 98.40\% ที่ FPR 0.00065 และเป็นอันดับหนึ่งในทั้งสามเกณฑ์การจัดอันดับ
    \item \textbf{ค่าการส่งผ่านสูงสุด} ภายใต้การแบ่งตามลำดับการจับข้อมูล ทิศ 2017\dirarrow{}2018 ได้สูงสุด 0.3368 (XGBoost) และทิศ 2018\dirarrow{}2017 ได้สูงสุด 0.4435 (LightGBM) เทียบกับค่าเพดานของแบบจำลองเดียวกัน 0.8589 และ 0.8767 และค่าพื้นประมาณ 0.13
    \item \textbf{ค่าดิบใกล้ค่าพื้น} ทิศ 2017\dirarrow{}2018 มี 4 จาก 7 แบบจำลองที่ได้เพียง 0.15--0.18 (ค่าพื้น 0.1337)
    \item \textbf{ความไม่สมมาตร} ทิศ 2018\dirarrow{}2017 กู้คืนได้มากกว่า โดยค่า Recovery เฉลี่ย 28.2\% เทียบกับ 14.4\% (ต่างกัน 13.8 จุด, 6 จาก 7 แบบจำลอง) ภายใต้การแบ่งตามลำดับการจับข้อมูล และ 15.2\% เทียบกับ 10.6\% (4.6 จุด, 4 จาก 7 แบบจำลอง) ภายใต้การแบ่งแบบสุ่ม
    \item \textbf{การแบ่งแบบสุ่มให้ค่าเพดานของ CICIDS2017 สูงกว่าการแบ่งตามลำดับการจับข้อมูล} ค่าเพดานเฉลี่ย 0.9065 (สุ่ม) เทียบกับ 0.8073 (ตามลำดับการจับข้อมูล) ขณะที่ CSE-CIC-IDS2018 ต่างกันเพียง 0.012
    \item \textbf{ความไวต่อสัดส่วนการแบ่ง} SSH-Patator ในชุดฝึกของ CICIDS2017 เพิ่มจาก 0 เป็น 474 ระเบียนเมื่อเปลี่ยนจาก 60\% เป็น 70\% ทำให้ความไม่สมมาตรลดจาก 33.6 เป็น 13.8 จุด และค่าเพดานเฉลี่ยของ CICIDS2017 เพิ่มจาก 0.649 เป็น 0.807
    \item \textbf{Brute Force} ค่า F1 ในทิศ 2017\dirarrow{}2018 เพิ่มจาก 0 เป็น 0.77 (LightGBM), 0.79 (XGBoost) และ 0.93 (Stacking) และอธิบายการเพิ่มของ Macro F1-score เฉลี่ยได้ 82\% (ที่เหลือ 18\% มาจาก Bot)
    \item \textbf{คลาสที่ส่งผ่านไม่ได้} DDoS ในทิศ 2017\dirarrow{}2018 ได้สูงสุดเพียง 0.011 และ Web Attack ได้สูงสุด 0.002 และ 0.018 ในสองทิศ ทั้งที่ค่าเพดานเฉลี่ยอยู่ที่ 0.99 และ 0.78--0.80
    \item \textbf{ความแตกต่างเฉพาะชุดข้อมูล} AUC ของการแยกคอร์ปัส 0.9990--1.0000 ทุกคลาส การสลับป้ายกำกับได้ 0.49--0.51 ภายในคอร์ปัสเดียวกันได้ 0.66--0.94 Logistic Regression ได้ตั้งแต่ 0.948 และไม่มีระเบียนซ้ำคร่อมชุดฝึกและชุดทดสอบ
    \item \textbf{SHAP} การเทียบภายในคอร์ปัสปลายทางผ่าน 0 จาก 11 กรณี (\textit{p} ต่ำสุด 0.174) และ \textit{G}\textsubscript{within} ของ DoS ทิศ 2017\dirarrow{}2018 ติดลบในทั้ง 4 แบบจำลอง
\end{enumerate}
